\documentclass[conference]{IEEEtran}

\usepackage[utf8]{inputenc}
\usepackage[indonesian]{babel}
\usepackage{amsmath,amssymb,amsfonts}
\usepackage{graphicx}
\usepackage{booktabs}
\usepackage{cite}
\usepackage[hidelinks]{hyperref}

\begin{document}

\title{Judul Proyek Akhir Peramalan Deret Waktu Sistem Tenaga Listrik\\
{\footnotesize \textsuperscript{*}Laporan Proyek Akhir Mata Kuliah Analisis Deret Waktu dan Peramalan (TEE-412)}}

\author{
\IEEEauthorblockN{1\textsuperscript{st} Nama Mahasiswa Pertama}
\IEEEauthorblockA{\textit{Program Studi S1 Teknik Elektro}\\
\textit{Universitas Bengkulu}\\
Bengkulu, Indonesia\\
npm1@unib.ac.id}
\and
\IEEEauthorblockN{2\textsuperscript{nd} Nama Mahasiswa Kedua}
\IEEEauthorblockA{\textit{Program Studi S1 Teknik Elektro}\\
\textit{Universitas Bengkulu}\\
Bengkulu, Indonesia\\
npm2@unib.ac.id}
\and
\IEEEauthorblockN{3\textsuperscript{rd} Nama Mahasiswa Ketiga}
\IEEEauthorblockA{\textit{Program Studi S1 Teknik Elektro}\\
\textit{Universitas Bengkulu}\\
Bengkulu, Indonesia\\
npm3@unib.ac.id}
}

\maketitle

\begin{abstract}
Tuliskan abstrak penelitian proyek akhir Anda secara ringkas ($\le 200$ kata) yang mencakup: (1) latar belakang permasalahan ketenagalistrikan (mis. peramalan beban jangka pendek, daya PLTS/PLTB, atau deteksi pencurian listrik \textit{smart meter}), (2) dataset dan protokol pemisahan kronologis bebas \textit{data leakage}, (3) model pembanding yang dievaluasi (minimal 1 \textit{baseline} statistik vs. 2 model ML/DL), serta (4) temuan kuantitatif utama (MAE, RMSE, MAPE/WAPE, atau F1-Score) dan implikasi keteknikannya.
\end{abstract}

\begin{IEEEkeywords}
analisis deret waktu, peramalan beban listrik, SARIMAX, XGBoost, LSTM, smart grid
\end{IEEEkeywords}

\section{Pendahuluan}
Jelaskan konteks rekayasa sistem tenaga yang melatarbelakangi studi kasus tim Anda, misalnya pentingnya akurasi \textit{Short-Term Load Forecasting} (STLF) atau kuantisasi ketidakpastian daya surya untuk penjadwalan \textit{Unit Commitment} dan BESS \cite{box2015, hyndman2021}. Rumuskan masalah, tujuan penelitian, serta ringkasan kontribusi tim Anda pada bagian akhir pendahuluan.

\section{Deskripsi Dataset \& Prapemrosesan}
Uraikan secara transparan sumber data, resolusi pencuplikan (mis. per-jam atau per-15 menit), periode pengamatan, dan statistik deskriptif utama. Jelaskan langkah prapemrosesan yang diterapkan:
\begin{itemize}
    \item Penanganan \textit{missing values} dan \textit{outlier} sensor.
    \item Analisis eksploratif: dekomposisi tren/musiman, uji stasioneritas ADF/KPSS, dan korelogram ACF/PACF.
    \item \textbf{Protokol Anti-\textit{Data Leakage}:} Pemisahan data latih, validasi, dan uji dilakukan secara kronologis (\textit{temporal split}). Seluruh transformasi skala (\texttt{MinMaxScaler}/\texttt{StandardScaler}) dan fitur \textit{rolling} dihitung secara kausal hanya dari histori masa lalu.
\end{itemize}

\section{Metodologi \& Arsitektur Model}
Bandingkan minimal satu model \textit{baseline} parametrik terhadap minimal dua model pembelajaran mesin / \textit{deep learning}. Tuliskan persamaan matematis inti, misalnya model SARIMA:
\begin{equation}
\Phi_P(B^s)\phi_p(B)(1-B)^d(1-B^s)^D y_t = \Theta_Q(B^s)\theta_q(B)\epsilon_t
\label{eq:sarima}
\end{equation}
serta metrik evaluasi utama seperti RMSE dan WAPE:
\begin{equation}
\text{RMSE} = \sqrt{\frac{1}{N}\sum_{t=1}^{N}(y_t - \hat{y}_t)^2}, \quad
\text{WAPE} = \frac{\sum_{t=1}^{N}|y_t - \hat{y}_t|}{\sum_{t=1}^{N}|y_t|} \times 100\%.
\label{eq:metrics}
\end{equation}
Jelaskan pula konfigurasi hiperparameter, perangkat komputasi (Google Colab/lokal), dan nilai \textit{random seed} untuk reproduksibilitas.

\section{Hasil Eksperimen \& Pembahasan}
Sajikan hasil evaluasi kuantitatif pada data uji (\textit{out-of-sample}) menggunakan tabel standar IEEE seperti pada Tabel~\ref{tab:hasil}. Lengkapi dengan grafik perbandingan sinyal aktual vs. prediksi, analisis residual (uji Ljung-Box / distribusi error), serta pembahasan kondisi ekstrem (hari libur, awan tebal, atau lonjakan beban puncak).

\begin{table}[htbp]
\caption{Contoh Perbandingan Kinerja Model Peramalan pada Data Uji}
\begin{center}
\begin{tabular}{lccc}
\toprule
\textbf{Model} & \textbf{MAE (MW)} & \textbf{RMSE (MW)} & \textbf{WAPE (\%)} \\
\midrule
Holt-Winters (Baseline) & 9.45 & 12.80 & 2.41 \\
SARIMAX (Cuaca) & 7.12 & 9.64 & 1.82 \\
XGBoost (Lag + Siklis) & 5.84 & 7.95 & 1.49 \\
LSTM Multi-Step (24 Jam) & \textbf{5.31} & \textbf{7.20} & \textbf{1.35} \\
\bottomrule
\end{tabular}
\label{tab:hasil}
\end{center}
\end{table}

Diskusikan implikasi praktis hasil tersebut terhadap operasi sistem tenaga listrik, misalnya pengurangan kebutuhan cadangan putar (\textit{spinning reserve}) atau penghematan biaya dispatch baterai BESS.

\section{Kesimpulan \& Saran}
Ringkas temuan paling signifikan yang menjawab rumusan masalah (didukung angka kinerja utama), sebutkan keterbatasan studi (mis. ketergantungan pada akurasi prakiraan cuaca), dan berikan saran pengembangan lanjutan.

\section*{Deklarasi Penggunaan AI \& Repositori Kode}
Tim menyatakan bahwa seluruh eksperimen dijalankan secara mandiri dan bebas \textit{data leakage}. Kode sumber lengkap (\texttt{.ipynb}, \texttt{requirements.txt}, dan \texttt{README.md}) tersedia pada: \url{https://github.com/username/nama-repositori-tim}. \textit{(Tuliskan secara jujur apabila menggunakan alat bantu AI untuk debugging atau penyuntingan bahasa.)}

\begin{thebibliography}{9}
\bibitem{box2015} G. E. P. Box, G. M. Jenkins, G. C. Reinsel, and G. M. Ljung, \textit{Time Series Analysis: Forecasting and Control}, 5th ed. Hoboken, NJ, USA: Wiley, 2015.
\bibitem{hyndman2021} R. J. Hyndman and G. Athanasopoulos, \textit{Forecasting: Principles and Practice}, 3rd ed. Melbourne, Australia: OTexts, 2021.
\bibitem{fei2020} S. Fei and Y. He, \textit{Load Forecasting in Electric Power Systems: Models and Methods}. Singapore: Springer, 2020.
\bibitem{benidis2022} K. Benidis et al., ``Deep learning for time series forecasting: Tutorial and literature survey,'' \textit{ACM Comput. Surv.}, vol. 55, no. 6, pp. 1--36, 2022.
\end{thebibliography}

\end{document}
